\begin{center}
    {\LARGE \textbf{2010分析化学}} \\[2ex]
    {\large 时量180分钟 \quad 满分90分}
\end{center}

\vspace{0.5cm}
\noindent\textbf{注意事项:}
\begin{enumerate}[label=\arabic*., parsep=0pt, itemsep=0pt]
    \item 答题前,考生先将自己的姓名、学号、考场号、座位号填写在答题卡上,将条形码准确粘贴在考生信息条形码粘贴区。
    \item 考试结束后,将本试卷与答题纸一并收回。
    \item 回答各个问题时,将答案写在答题纸上,写在本试卷上无效。
\end{enumerate}

\vspace{0.5cm}
\noindent\textbf{一、简答题(28 pts.)}

\begin{enumerate}[label=\arabic*., itemsep=1.5ex]
    \item 写出下列溶液的质子条件式
    \begin{enumerate}[label=(\arabic*), noitemsep, topsep=0pt, partopsep=0pt]
        \item $0.10\ \mathrm{mol/L}\ \ce{HCl}$溶液和$0.10\ \mathrm{mol/L}\ \ce{H2SO4}$混合所得溶液；
        \item $\ce{Na2CO3}$溶液
    \end{enumerate}
    \item 在弱碱性溶液中用EDTA滴定$\ce{Zn^{2+}}$常使用$\ce{NH3-NH4+}$溶液，其作用是什么？
    \item (1) 指出下列数字有效数字的位数：$\mathrm{pH}=10.46$，$c=0.02420\ \mathrm{mol/L}$
    \\
    (2) 按有效数字修约规则将$2.45651$和$2.4565$修约为四位有效数字，分别是多少？
    \item 测定$\ce{(NH4)2SO4}$中的氮时，能否用$\ce{NaOH}$标准溶液直接滴定？为什么？（$\ce{NH3}$的$\mathrm{p}K_\mathrm{b}=4.74$）
    \item 符合朗伯-比尔定律的一有色溶液，当有色物质的浓度增加时，最大吸收波长和吸光度各有何变化？
    \item 写出两种标定盐酸溶液的基准物质。
    \item 在碘量法测定铜合金中铜的含量的实验中，加入$\ce{KSCN}$的作用是什么？$\ce{KSCN}$不能过早加入，为什么？
\end{enumerate}

\vspace{0.5cm}
\noindent\textbf{二、计算题(62 pts. total)}

\begin{enumerate}[label=\arabic*., itemsep=1.5ex]
    \item (8 pts.) 丁酸在苯-水体系中的分配系数$K_D=3$，当$100\ \mathrm{mL}\ 0.10\ \mathrm{mol/L}$丁酸溶液用$25\ \mathrm{mL}$苯萃取时，计算在$\mathrm{pH}=10.00$时丁酸在苯中的浓度？（丁酸根离子不被萃取、丁酸的$\mathrm{p}K_\mathrm{a}=4.8$）
    \item (10 pts.) 已知在$1\ \mathrm{mol/L}\ \ce{HCl}$溶液中滴定$E^{\ominus\prime}(\ce{Fe^{3+}/Fe^{2+}})=0.68\ \mathrm{V}$，$E^{\ominus\prime}(\ce{Sn^{4+}/Sn^{2+}})=0.14\ \mathrm{V}$
    \begin{enumerate}[label=(\arabic*), noitemsep, topsep=0pt, partopsep=0pt]
        \item 若用$\ce{Fe^{3+}}$的$\ce{HCl}$溶液滴定$\ce{SnCl2}$溶液，计算达到化学计量点时的电位是多少？
        \item 若将$0.10\ \mathrm{mol/L}\ \ce{FeCl3}$溶液与$0.10\ \mathrm{mol/L}\ \ce{SnCl2}$溶液等体积混合，平衡时体系电位是多少？
    \end{enumerate}
    \item (10 pts.) 电分析法测定某患者血糖的浓度（$\mathrm{mmol/L}$），10次结果为：$7.5$，$7.4$，$7.7$，$7.6$，$7.5$，$7.6$，$7.6$，$7.5$，$7.6$，$7.6$，求相对标准差及置信度$95\%$的置信区间，此结果与正常人血糖的质量分数$6.7\ \mathrm{mmol/L}$是否有显著性差异？
    \\
    已知：
    \begin{center}
    \begin{tabular}{l|ccc}
    $f$ & $8$ & $9$ & $10$ \\
    \hline
    $t_{0.05}$ & $2.31$ & $2.26$ & $2.23$ \\
    \end{tabular}
    \end{center}
    \item (12 pts.) 用$0.02\ \mathrm{mol/L}\ \ce{NaOH}$溶液滴定$0.02\ \mathrm{mol/L}\ \ce{HCl}$和$0.02\ \mathrm{mol/L}\ \ce{HAc}$的混合溶液，能否选择滴定$\ce{HCl}$和连续滴定$\ce{HAc}$；滴至$\mathrm{pH}=9.0$时的终点误差是多少？（$\mathrm{p}K_\mathrm{a}(\ce{HAc})=4.74$）
    \item (12 pts.) 用$2\times10^{-2}\ \mathrm{mol/L}\ \mathrm{EDTA}$滴定同浓度$\ce{Zn^{2+}}$，$\ce{Al^{3+}}$混合液中的$\ce{Zn^{2+}}$，以磺基水杨酸($\ce{H2L}$)掩蔽$\ce{Al^{3+}}$。若终点时$\mathrm{pH}$为$5.5$，选二甲酚橙为指示剂。计算终点误差。（已知$\lg K(\ce{ZnY})=16.5$；$\lg K(\ce{AlY})=16.1$；$\mathrm{pH}$为$5.5$时，$\lg\alpha_{Y(\mathrm{H})}=5.5$，$\mathrm{pZn_{(二甲酚橙)}}=5.7$；$\lg\alpha_{Al(\mathrm{L})}=6.4$）
    \item (10 pts.) 在$\mathrm{pH}=2.0$的含有$0.01\ \mathrm{mol/L}\ \mathrm{EDTA}$及$0.10\ \mathrm{mol/L}\ \ce{HF}$的溶液中，当加入$\ce{CaCl2}$使溶液中的$c(\ce{Ca^{2+}})=0.10\ \mathrm{mol/L}$时，问
    \begin{enumerate}[label=(\arabic*), noitemsep, topsep=0pt, partopsep=0pt]
        \item EDTA的存在对生成$\ce{CaF2}$沉淀有无影响？
        \item 能否产生$\ce{CaF2}$沉淀？(不考虑体积变化) （已知$\ce{HF}$的$K_\mathrm{a}=6.6\times10^{-4}$，$K_{\mathrm{sp}}(\ce{CaF2})=2.7\times10^{-11}$，$\lg K(\ce{CaY})=10.69$，$\lg\alpha_{Y(\mathrm{H})}=13.51$）
    \end{enumerate}
\end{enumerate}